% ============================================================
% CHAPITRE 2
% ============================================================
\chapter{Blanchiment et analyse de graphes : état de l'art}
\label{ch:etatart}

\section{Les trois temps du blanchiment}
Le blanchiment consiste à donner une origine apparemment licite à des fonds
qui n'en ont pas. On le décrit classiquement en trois étapes
(figure~\ref{fig:etapes}). Les typologies visées par ce projet se situent
surtout dans la deuxième, l'\emph{empilement} : multiplier les transferts
pour brouiller la trace. C'est précisément l'étape qui laisse une empreinte
dans la \emph{structure} du réseau de virements.

\begin{figure}[htp]
\centering
\begin{tikzpicture}[
  etape/.style={rounded corners=3pt, minimum height=1.1cm, text width=4.3cm, align=center,
    font=\sffamily\small\bfseries, text=white, fill=#1},
  detail/.style={font=\sffamily\scriptsize, text width=4.3cm, align=left, anchor=north},
  typo/.style={rounded corners=2pt, fill=#1!15, draw=#1, font=\sffamily\tiny\bfseries, inner sep=2.5pt}]
  \node[etape=grismoyen] (p) at (0,0) {1 · Placement};
  \node[etape=orangeprojet] (e) at (5.2,0) {2 · Empilement (\emph{layering})};
  \node[etape=bleuprojet] (i) at (10.4,0) {3 · Intégration};
  \draw[fleche] (p) -- (e);
  \draw[fleche] (e) -- (i);
  \node[detail] at (0,-0.75) {Faire entrer l'argent dans le système financier, en petites sommes pour
    rester sous les seuils de déclaration.};
  \node[detail] at (5.2,-0.75) {Multiplier les virements entre comptes et sociétés écrans, faire
    tourner l'argent, le fractionner puis le regrouper.};
  \node[detail] at (10.4,-0.75) {Réinjecter les fonds dans l'économie légale : achats, investissements,
    factures de complaisance.};
  \node[typo=cat5] at (0,-2.65) {mules (smurfing)};
  \node[typo=cat1] at (3.75,-2.65) {carrousel};
  \node[typo=cat3] at (5.25,-2.65) {cycle long};
  \node[typo=cat4] at (6.95,-2.65) {scatter-gather};
  \node[typo=cat2] at (4.45,-3.15) {pivot};
  \node[typo=cat6] at (6.05,-3.15) {réseau lié};
  \node[font=\sffamily\tiny, color=encre2] at (10.4,-2.65) {hors du périmètre du flux de virements};
  \node[font=\sffamily\tiny\bfseries, color=encre2, anchor=east] at (-2.3,-2.9) {alertes :};
\end{tikzpicture}
\caption[Les trois étapes du blanchiment et les typologies détectées]{Les trois étapes du blanchiment
et la place des six types d'alertes produits par la plateforme.}
\label{fig:etapes}
\end{figure}

\section{Cinq typologies, cinq formes de graphe}
Les typologies simulées reprennent celles du jeu de données synthétique
d'IBM pour l'anti-blanchiment~\cite{ibmaml} (\emph{fan-in}, cycle,
\emph{scatter-gather}), auxquelles s'ajoute le compte pivot qui relie
plusieurs carrousels. La figure~\ref{fig:typologies} les dessine telles que le
producteur les injecte ; chacune appelle un algorithme de graphe différent.

\begin{figure}[htp]
\centering
\begin{tikzpicture}[
  v/.style={circle, fill=#1, minimum size=4.6mm, inner sep=0pt},
  e/.style={-{Stealth[length=1.5mm]}, line width=0.7pt, color=encre2},
  titre/.style={font=\sffamily\scriptsize\bfseries, anchor=north},
  sous/.style={font=\sffamily\tiny, color=encre2, align=center, text width=2.7cm, anchor=north}]
  % 1. carrousel
  \begin{scope}[xshift=0cm]
    \node[v=cat1] (a) at (90:0.75) {}; \node[v=cat1] (b) at (-30:0.75) {}; \node[v=cat1] (c) at (210:0.75) {};
    \draw[e, bend left=20] (a) to (b); \draw[e, bend left=20] (b) to (c); \draw[e, bend left=20] (c) to (a);
    \node[titre] at (0,-1.05) {Carrousel};
    \node[sous] at (0,-1.4) {3 sociétés, 40--250\,k MAD, $-$1 à 3\,\%/étape, 1--3 min\\[2pt]\textbf{motif à 3 arêtes}};
  \end{scope}
  % 2. cycle long
  \begin{scope}[xshift=3.15cm]
    \foreach \k in {0,...,4} \node[v=cat3] (n\k) at (90+\k*72:0.78) {};
    \foreach \k [evaluate=\k as \j using {int(mod(\k+1,5))}] in {0,...,4} \draw[e] (n\k) -- (n\j);
    \node[titre] at (0,-1.05) {Cycle long};
    \node[sous] at (0,-1.4) {4 ou 5 comptes, même signature que le carrousel\\[2pt]\textbf{SCC + motif à 4--5 arêtes}};
  \end{scope}
  % 3. mule
  \begin{scope}[xshift=6.3cm]
    \node[v=cat5, minimum size=6mm] (m) at (0,0) {};
    \foreach \k in {0,...,7} { \node[v=grismoyen, minimum size=3.2mm] (s\k) at (100+\k*22:0.95) {}; \draw[e] (s\k) -- (m); }
    \node[v=grismoyen, minimum size=3.6mm] (bf) at (-30:0.95) {};
    \draw[e, line width=1.4pt, color=cat5] (m) -- (bf);
    \node[titre] at (0,-1.05) {Compte mule};
    \node[sous] at (0,-1.4) {8--14 dépôts de 2--10\,k MAD, puis un retrait de 90--97\,\%\\[2pt]\textbf{vélocité + flux traversant}};
  \end{scope}
  % 4. scatter-gather
  \begin{scope}[xshift=9.45cm]
    \node[v=cat4] (s) at (-0.95,0) {}; \node[v=cat4] (t) at (0.95,0) {};
    \foreach \k in {-2,...,2} { \node[v=grismoyen, minimum size=3.6mm] (m\k) at (0,\k*0.33) {};
      \draw[e] (s) -- (m\k); \draw[e] (m\k) -- (t); }
    \node[titre] at (0,-1.05) {Scatter-gather};
    \node[sous] at (0,-1.4) {4--6 intermédiaires, parts égales, 97--99\,\% reversés\\[2pt]\textbf{motif à 2 sauts}};
  \end{scope}
  % 5. pivot
  \begin{scope}[xshift=12.6cm]
    \node[v=cat2, minimum size=6mm] (p) at (0,0) {};
    \foreach \k in {0,1,2} {
      \node[v=grismoyen, minimum size=3.6mm] (x\k) at (60+\k*120:0.95) {};
      \node[v=grismoyen, minimum size=3.6mm] (y\k) at (100+\k*120:0.95) {};
      \draw[e] (p) -- (x\k); \draw[e] (x\k) -- (y\k); \draw[e] (y\k) -- (p); }
    \node[titre] at (0,-1.05) {Pivot};
    \node[sous] at (0,-1.4) {1 compte au centre de 3 carrousels imbriqués\\[2pt]\textbf{cycles + PageRank, puis CC}};
  \end{scope}
\end{tikzpicture}
\caption[Les cinq typologies injectées et leur forme dans le graphe]{Les cinq typologies injectées,
leurs paramètres de génération et l'algorithme qui les détecte (en gras).}
\label{fig:typologies}
\end{figure}

\section{Les algorithmes de graphe mobilisés}
Le tableau~\ref{tab:algos} résume ce que chaque algorithme sait dire d'un
graphe de virements, et ce qu'il coûte. Le point clé est l'ordre de grandeur :
la recherche de motif est une suite de jointures dont le coût explose avec la
taille du graphe ; PageRank et les composantes sont itératifs et ne
s'exécutent raisonnablement que sur un graphe réduit. Toute l'architecture du
chapitre suivant découle de ce constat.

\begin{table}[htp]
\centering\small
\begin{tabular}{@{}>{\raggedright\arraybackslash}p{2.7cm}>{\raggedright\arraybackslash}p{4.4cm}>{\raggedright\arraybackslash}p{3.7cm}>{\raggedright\arraybackslash}p{3.3cm}@{}}
\toprule
\sffamily\bfseries Algorithme & \sffamily\bfseries Question posée & \sffamily\bfseries Coût & \sffamily\bfseries Usage ici\\
\midrule
Motif finding & Existe-t-il un sous-graphe de cette forme ? & 1 jointure par arête du motif ; explose avec le degré des hubs & carrousels (rapide), scatter-gather, cycles 4--5\\
PageRank & Quels comptes concentrent les flux ? & itératif (10 tours ici) & comptes pivots\\
SCC & Quels comptes peuvent tous s'atteindre en suivant le sens des virements ? & itératif ; inutile hors d'un noyau élagué & cycles longs\\
Connected Components & Quels comptes sont reliés, quel que soit le sens ? & itératif & réseaux d'alertes liées\\
Degré, vélocité & Combien de dépôts, à quel rythme ? & linéaire & comptes mules\\
\bottomrule
\end{tabular}
\caption{Algorithmes de graphe utilisés, coût et usage.}
\label{tab:algos}
\end{table}

\section{Approches de détection}
Les dispositifs bancaires reposent historiquement sur des règles par
transaction (seuils, listes) puis sur des modèles d'apprentissage notant
chaque transaction ou chaque client. Aucune de ces deux familles ne voit une
boucle : elle n'est la propriété d'aucun élément isolé. Les analyses de
graphe la voient, mais sont souvent conduites en différé, sur des
historiques. Le projet se place dans la quatrième case du
tableau~\ref{tab:approches} : l'analyse de graphe \emph{en flux}.

\begin{table}[htp]
\centering\small
\begin{tabular}{@{}>{\raggedright\arraybackslash}p{4.0cm}cccc@{}}
\toprule
& \sffamily\bfseries Règles & \sffamily\bfseries ML tabulaire & \sffamily\bfseries Graphe (différé) & \sffamily\bfseries Graphe en flux\\
\midrule
Voit les boucles et les réseaux & \non & \non & \oui & \oui\\
Alerte pendant que l'argent circule & \oui & \oui & \non & \oui\\
Explication lisible par un analyste & \oui & \partiel & \partiel & \oui\\
Se passe de données étiquetées & \oui & \non & \oui & \oui\\
Simplicité de mise en œuvre & \oui & \partiel & \partiel & \non\\
\bottomrule
\end{tabular}
\caption[Positionnement parmi les approches de détection]{Positionnement parmi les approches de
détection (\oui~oui, \partiel~en partie, \non~non). Le prix de la dernière colonne est la complexité
technique, objet des chapitres~\ref{ch:architecture} et~\ref{ch:performance}.}
\label{tab:approches}
\end{table}

\section{Cadre marocain}
Au Maroc, la lutte contre le blanchiment repose sur la loi n\up{o}\,43-05,
renforcée par la loi n\up{o}\,12-18 qui a notamment institué l'Autorité
Nationale du Renseignement Financier (ANRF), destinataire des déclarations de
soupçon des banques~\cite{anrf}. Le pays a été retiré en 2023 de la liste des
juridictions sous surveillance renforcée du GAFI~\cite{gafi}. Pour une banque,
l'enjeu est double : déclarer vite, et pouvoir justifier chaque déclaration.
C'est ce que vise la plateforme en associant à chaque alerte un score
\emph{expliqué}, son sous-graphe et un historique non modifiable.

\begin{figure}[htp]
\centering
\begin{tikzpicture}[x=1.1cm]
  \draw[grismoyen, line width=0.8pt] (0,0) -- (11.6,0);
  \foreach \x/\an/\txt/\c in {
     0.9/2007/{Loi 43-05 : incrimination du blanchiment, obligations de vigilance}/bleuprojet,
     4.0/2009/{Cellule de renseignement financier opérationnelle}/grismoyen,
     7.2/2021/{Loi 12-18 : renforcement, création de l'ANRF}/bleuprojet,
     10.3/2023/{Sortie de la liste grise du GAFI}/vertcode} {
    \fill[\c] (\x,0) circle (1.4mm);
    \node[font=\sffamily\small\bfseries, color=\c, anchor=south] at (\x,0.15) {\an};
    \node[font=\sffamily\tiny, text width=2.9cm, align=center, anchor=north] at (\x,-0.2) {\txt};
  }
  \node[font=\sffamily\tiny\bfseries, color=orangeprojet, anchor=west, text width=2.1cm] at (11.8,0.0)
    {ce projet : détecter et documenter en temps réel};
\end{tikzpicture}
\caption{Repères du cadre marocain de lutte contre le blanchiment.}
\label{fig:cadre}
\end{figure}
